\begin{initiativkarte}{Peer to Peer}

  % Bild (optional)
  \IfFileExists{bilder/peer_to_peer.jpg}{%
    \begin{center}
      \includegraphics[width=\linewidth,height=5cm,keepaspectratio]{bilder/peer_to_peer.jpg}
    \end{center}
    \vspace{0.4cm}
  }{}

  % Kurzbeschreibung
  \textit{\textcolor{hawgray}{Peer to Peer bietet Unterstützung für Studierende in belastenden Situationen – von Studierenden für Studierende. Wir helfen bei Stress, Überforderung und persönlichen Herausforderungen im Studium.}}

  \vspace{0.5cm}
  \hawrule

  % Beschreibung
  \textbf{\textcolor{hawblue}{Was machen wir?}}\\[0.3cm]

  Peer to Peer ist ein Unterstützungsangebot der HAW Hamburg, bestehend aus Studierenden verschiedener Fakultäten. Wir kennen die Herausforderungen des Studienalltags aus eigener Erfahrung und möchten anderen Studierenden in schwierigen Situationen helfen.

  Unser Fokus liegt auf der Unterstützung von Studierenden mit psychischen Belastungen, chronischen Erkrankungen oder generellen Schwierigkeiten im Studienalltag.

  Wir unterstützen dich unter anderem bei:
  \begin{itemize}
      \item Studienorganisation und Zeitmanagement
      \item Finanzierung des Studiums
      \item Urlaubssemester und Wiedereinstieg
      \item Nachteilsausgleichen
      \item Prüfungsstress und Überforderung
      \item persönlichen Gesprächen und emotionaler Unterstützung
  \end{itemize}

  Zusätzlich bieten wir:
  \begin{itemize}
      \item regelmäßige offene Sprechzeiten
      \item längerfristige Begleitung
      \item moderierte Gruppenangebote zum Austausch mit anderen Studierenden
  \end{itemize}

  \vspace{0.5cm}

  % Tags
  \tag{Beratung}
  \tag{Psychische Gesundheit}
  \tag{Unterstützung}
  \tag{Studium}
  \tag{Community}

  \vspace{0.6cm}
  \hawrule

  % Infozeilen
  \infozeile{\faUsers}{Zielgruppe: Alle Studierenden der HAW Hamburg}
  \infozeile{\faGraduationCap}{Semester: Ab dem 1. Semester}
  \infozeile{\faMapMarker*}{Ort: HAW Hamburg (verschiedene Standorte / Sprechzeiten)}
  \infozeile{\faEnvelope}{Kontakt: Über Linktree}
  \infozeile{\faGlobe}{Website: \href{https://linktr.ee/p2p_hawhamburg}{linktr.ee/p2p\_hawhamburg}}

\end{initiativkarte}